\scanpage{94}
% loc 1, page identifier
PU \mspage{\handwritten{\unsure{65}}}

% loc 2–6, upper crossed paragraph
\deleted{Cette abstraction est \deleted{apparente en}\inserted{\handwritten{inhérente à}} tout mot \deleted{sans exception} du langage\inserted{ sans exception}\inserted{\handwritten{\unsure{,|.}}} \deleted{\unsure{et}} \handwritten{C'}est \deleted{même}\handwritten{bien} la \inserted{\handwritten{aussi une}} condition \deleted{sine qua non}\handwritten{nécessaire}, pour que le mot soit bien le patrimoine culturel d'un groupe ayant une continuité dans le temps, et non la propriété d'une personne\inserted{\handwritten{ ni d'un groupement de personnes}}. Il en est ainsi même pour les mots les plus élémentaires\inserted{\handwritten{ de tous}}, comme ``mère'', ``père'', ``manger'', ``boire'', ``soleil'' etc.}
% loc 4, left margin
\margin{\handwritten{\unsure{d très|très}}}
% loc 7–12
\deleted{Entre de tels mots-\deleted{\illegible{1}}\handwritten{base}\deleted{, \unsure{en}} et des mots comme ``famille'', voire ``groupe'', ``peuple'', ``nation'', ``gouvernement'', ``politique'', ``philosophie'', ``abstraction'', on sent \deleted{\unsure{une} une différence qualitative, \unsure{marquée}, une différence de ``degré d'abstraction''}\inserted{\handwritten{une ``distance''}}, comparable à celle qui existe entre les \deleted{mots ``élémentaires''}\handwritten{``mots-base''} eux-mêmes, et les ``objets concrets'' \deleted{(au sens le plus strict\unsure{,} de tantôt)} auxquels ils s'appliquent.}

% loc 13–21, second paragraph
\deleted{Pour le dire autrement: le processus d'abstraction, qui \deleted{est}\inserted{semble être} (au point de vue de la formation de \uline{concepts}, porteurs de \uline{sens}) \deleted{l'opération essentielle}\handwritten{comme l'âme} dans la formation progressive d'un langage de plus en plus complexe, de plus en plus ``ramifié'', \deleted{et} apte à saisir \deleted{\unsure{les}}des replis et \deleted{\unsure{les}}des nuances \deleted{les} de plus en plus délicats\deleted{\unsure{e}} des choses de l'Univers -- ce processus est c\,u\,m\,u\,l\,a\,t\,i\,f . Une fois un de ses pas accomplis, le nouveau concept, s'incarnant dans un mot nouveau, finit par faire partie des choses familières\inserted{ et ``concrètes''}, au même titre, voire plus encore, que ces autres choses que nous avions qualifiées\inserted{\handwritten{ tantôt}} de \deleted{\unsure{``}}v\,r\,a\,i\,m\,e\,n\,t concrètes. \deleted{C'est sûr que pour un mathématicien}}
% loc 22–34, paragraph sign before Ainsi
\margin{\handwritten{\unsure{\P}}}
\deleted{Ainsi, le concept de n\,o\,m\,b\,r\,e (entier, disons), historiquement ou structuralement parlant, est\inserted{\handwritten{ visiblement}} d'un degré d'abstraction élevé (bien plus élevé qu'un mot-concept comme ``deux'' ou ``trois'' par exemple, \handwritten{le}qu\handwritten{el} représente\deleted{\unsure{nt}} \deleted{\unsure{une qualité|une qualités}} tout ce qu'il y a de terre-à-terre). Cela n'empêche qu'il n'y a pas un seul mathématicien, \inserted{sûrement,} à qui cette notion n'appara\handwritten{îsse} comme extrêmement \handwritten{``}concrète\handwritten{''} \deleted{par contre} -- bien plus concrète disons que ``feu'' ou ``terre'', \deleted{qui} désign\handwritten{an}t des choses avec lesquelles il a plus ou moins perdu contact (à supposer qu'il l'ait jamais eu). Au niveau d'abstraction où travaillent \deleted{la plupart}\inserted{beaucoup} des mathématiciens, le nombre \deleted{entier} apparaît aussi \inserted{\handwritten{concret (et en même temps, aussi}} lointain!\handwritten{)}, depuis les hauteurs où il \deleted{fait}\handwritten{trace} ses virevoltes, que \deleted{apparaîtra}\inserted{\handwritten{le}} serait le garde-manger de \handwritten{s}a maison familiale, à celui qui parlerait de l'économie française ou du marché mondial des céréales.}

% loc 35–40, final paragraph
\deleted{Tout ceci pour illustrer qu\deleted{e}\handwritten{'} \deleted{la notion d'``abstraction'',} à toutes fins pratiques et une fois réalisé que \uline{toute} pensée et \uline{tout} langage sont ``abstraits'' (par rapport aux ``objets'' au sens courant du terme), \inserted{\handwritten{la notion d'``abstraction'',}} est avant tout une \deleted{\illegible{1}} n\,o\,t\,i\,o\,n r\,e\,l\,a\,t\,i\,v\,e . Un discours, un langage, une pensée, une thé\unsure{orie|oire} sont ``plus abstraits'' ou ``plus concrets'' qu'un\inserted{\handwritten{ (ou une)}} autre. Quand l'un et l'autre\inserted{\handwritten{ (ou l'une et l'autre)}} appartiennent à des domaines pas trop éloignés de la}
% archive sheet number: 93
